\documentclass{article}

\usepackage{palatino,eulervm,dutchcal,xcolor}
\usepackage{graphicx,subcaption,float}
\usepackage{enumitem,parskip,multicol,array}
\newcolumntype{L}{>{$}c<{$}}%table in mathmode
\usepackage{amsthm,amssymb,amsmath,mathtools,thmtools,tikz,tikz-cd}
\usetikzlibrary{%
	matrix,%
	calc,%
	arrows,%
	shapes,
	decorations.markings,backgrounds,calc,intersections}
\tikzcdset{scale cd/.style={every label/.append style={scale=#1},
		cells={nodes={scale=#1}}}}
\usepackage[bookmarks,bookmarksopen,bookmarksdepth=3]{hyperref}
\hypersetup{%colores
	colorlinks=true,
	urlcolor=blue,
	linkcolor=magenta,
	citecolor=blue,
	filecolor=blue,
	urlbordercolor=white,
	linkbordercolor=white,
	citebordercolor=white,
	filebordercolor=white}
\usepackage{cleveref}

\definecolor{blue-violet}{rgb}{0.54, 0.17, 0.89}
\definecolor{azure}{rgb}{0.0, 0.5, 1.0}
\definecolor{green(ncs)}{rgb}{0.0, 0.62, 0.42}
\definecolor{forestgreen}{rgb}{0.13, 0.55, 0.13}
\definecolor{limegreen}{rgb}{0.2, 0.8, 0.2}
\definecolor{palatinateblue}{rgb}{0.15, 0.23, 0.89}
\definecolor{trueblue}{rgb}{0.0, 0.45, 0.81}
\definecolor{goldenyellow}{rgb}{1.0, 0.87, 0.0}
\definecolor{fashionfuchsia}{rgb}{0.96, 0.0, 0.63}
\definecolor{brightcerulean}{rgb}{0.11, 0.67, 0.84}
\definecolor{jonquil}{rgb}{0.98, 0.85, 0.37}
\definecolor{lavendermagenta}{rgb}{0.93, 0.51, 0.93}
\definecolor{peru}{rgb}{0.8, 0.52, 0.25}
\definecolor{persimmon}{rgb}{0.93, 0.35, 0.0}
\definecolor{persianred}{rgb}{0.8, 0.2, 0.2}
\definecolor{persianblue}{rgb}{0.11, 0.22, 0.73}
\definecolor{persiangreen}{rgb}{0.0, 0.65, 0.58}
\definecolor{persianyellow}{rgb}{0.9, 0.89, 0.0}


\declaretheoremstyle[headfont=\color{trueblue}\normalfont\bfseries,]{colored1}
\declaretheoremstyle[headfont=\color{forestgreen}\normalfont\bfseries,]{colored2}
\declaretheoremstyle[headfont=\color{peru}\normalfont\bfseries,]{colored3}
\declaretheoremstyle[headfont=\color{persiangreen}\normalfont\bfseries,]{colored4}
\declaretheoremstyle[headfont=\color{brightcerulean}\normalfont\bfseries,]{colored5}
\declaretheoremstyle[headfont=\color{lavendermagenta}\normalfont\bfseries,]{colored6}
\declaretheoremstyle[headfont=\color{blue-violet}\normalfont\bfseries,]{colored7}
\declaretheoremstyle[headfont=\color{green(ncs)}\normalfont\bfseries,]{colored8}
\declaretheoremstyle[headfont=\color{peru}\normalfont\bfseries,]{colored9}
\declaretheoremstyle[headfont=\color{persiangreen}\normalfont\bfseries,]{colored10}

\declaretheorem[style=colored1,name=Theorem]{thm}
\declaretheorem[style=colored2,numberlike=thm,name=proposition]{prop}
\declaretheorem[style=colored3,numbered=no,name=Lemma]{lemma}
\declaretheorem[style=colored4,numbered=no,name=Corollary]{coro}
\declaretheorem[style=colored5,numbered=no,name=Example]{example}
\declaretheorem[style=colored5,numbered=no,name=Examples]{exemplos}
\declaretheorem[style=colored7,numbered=no,name=Exercise]{exercise}
\declaretheorem[style=colored6,numbered=no,name=Remark]{remark}
\declaretheorem[style=colored8,numbered=no,name=Claim]{claim}
\declaretheorem[style=colored9,numbered=no,name=Definition]{definition}
\declaretheorem[style=colored10,numbered=no,name=Question]{question}

\numberwithin{equation}{section}

\newcommand{\A}{\mathbb{A}}
\newcommand{\R}{\mathbb{R}}
\newcommand{\Z}{\mathbb{Z}}
\newcommand{\N}{\mathbb{N}}
\newcommand{\C}{\mathbb{C}}
\newcommand{\Q}{\mathbb{Q}}
\newcommand{\D}{\mathbb{D}}
\renewcommand{\P}{\mathbb{P}}

\newcommand{\Ac}{\mathcal{A}}
\newcommand{\Bc}{\mathcal{B}}
\newcommand{\Cc}{\mathcal{C}}
\newcommand{\Dc}{\mathcal{D}}
\newcommand{\Ec}{\mathcal{E}}
\newcommand{\Fc}{\mathcal{F}}
\newcommand{\Gc}{\mathcal{G}}
\newcommand{\Hc}{\mathcal{H}}
\newcommand{\Ic}{\mathcal{I}}
\newcommand{\Jc}{\mathcal{J}}
\newcommand{\Kc}{\mathcal{K}}
\newcommand{\Lc}{\mathcal{L}}
\newcommand{\Mc}{\mathcal{M}}
\newcommand{\Nc}{\mathcal{N}}
\newcommand{\Oc}{\mathcal{O}}
\newcommand{\Pc}{\mathcal{P}}
\newcommand{\Qc}{\mathcal{Q}}
\newcommand{\Rc}{\mathcal{R}}
\newcommand{\Sc}{\mathcal{S}}
\newcommand{\Tc}{\mathcal{T}}
\newcommand{\Uc}{\mathcal{U}}
\newcommand{\Vc}{\mathcal{V}}
\newcommand{\Wc}{\mathcal{W}}
\newcommand{\Xc}{\mathcal{X}}
\newcommand{\Yc}{\mathcal{Y}}
\newcommand{\Zc}{\mathcal{Z}}

\DeclareMathOperator{\img}{img}
\DeclareMathOperator{\Arg}{Arg}
\DeclareMathOperator{\id}{id}
\DeclareMathOperator{\pt}{pt}
\DeclareMathOperator{\Alt}{Alt}
\DeclareMathOperator{\sgn}{sgn}
\DeclareMathOperator{\hTop}{h-Top}
\DeclareMathOperator{\supp}{supp}
\DeclareMathOperator{\Int}{Int}
\DeclareMathOperator{\Ob}{Ob}
\DeclareMathOperator{\Mor}{Mor}
\DeclareMathOperator{\Top}{Top}
\DeclareMathOperator{\Set}{Set}
\DeclareMathOperator{\CGWH}{CGWH}
\DeclareMathOperator{\Hom}{Hom}
\DeclareMathOperator{\Map}{Map}
\DeclareMathOperator{\Tot}{Tot}
\DeclareMathOperator{\op}{op}
\DeclareMathOperator{\ev}{ev}
\DeclareMathOperator{\hofib}{hofib}
\DeclareMathOperator{\rel}{rel}
\DeclareMathOperator{\Nat}{Nat}
\DeclareMathOperator{\Arr}{Arr}
\DeclareMathOperator{\Funct}{Funct}
\DeclareMathOperator{\colim}{colim}
\DeclareMathOperator{\GL}{GL}
\DeclareMathOperator{\SL}{SL}
\DeclareMathOperator{\SO}{SO}
\DeclareMathOperator{\U}{U}
\DeclareMathOperator{\SU}{SU}
\DeclareMathOperator{\Sp}{Sp}
\DeclareMathOperator{\M}{M}
\DeclareMathOperator{\Aut}{Aut}
\DeclareMathOperator{\PGL}{PGL}
\DeclareMathOperator{\PSL}{PSL}
\DeclareMathOperator{\Pic}{Pic}
\DeclareMathOperator{\Loc}{Loc}
\DeclareMathOperator{\Alb}{Alb}
\DeclareMathOperator{\rk}{rk}

\begin{document}
	{\Large Class Field Theory (CFT) and geometric CFT}

	Sources
\begin{enumerate}
	\item B. Conrad. Geometric CFT.
	\item E. Frenkel. Geometric Langlands and CFT.
	\item B. Bhatt. Towards geometric CFT.
	\item G. Vooys. Local Langlands for $\GL(1)$.
\end{enumerate}

Part 1
\begin{enumerate}
	\item The canonical extension (after Deligne, exposition by Bhott.
	\item 1. $\implies$ unramified global CFT.
	\item CFT main results (after \ldots exposition due to Yihan Zhu)
	\item Global geometric CFT (exposition by Conrad)
	\item Local geometric CFT (after by Ronsenficht \& Long, exposition due to Vooys'20)
\end{enumerate}

Let $X/\C$ be a projective irreducible curve. Consider the map
\begin{align*}
	AJ:X  &\longrightarrow \Pic^1(X) \\
	x &\longmapsto (x) = [x]
\end{align*}
\begin{definition}
$\Loc_1(X)$ is the \textit{\textbf{category of $\rk=1$ local systems over X}}. 	
\end{definition}
This is the category of representations of the fundamental group, that is
\begin{align*}
	\pi_0(\Loc_1(X))&=\{\pi_1(X)\to \C^*\}\\
	\pi_0(\Loc_1(\Alb(X))&=\{H_1(X,\Z)\to \C^*\} 
\end{align*}
Some more comments

\section{some other things}

There's some algebraic geometry about all this. 

What's evaluation. Consider the completion of real numbers from rational numbers. This is completion with respect to the norm $|\cdot|$. But we may also complete respect to $|\cdot|_p$ for a prime number $p$. It goes like this. For any rational number $q$ 
	\begin{align*}
		\Q &\longrightarrow \Q_{>0} \\
		r &\longmapsto |r|\\
		r &\longmapsto |r|_p:=\text{the exponent of $p$ in the prime factorization}
	\end{align*}

	$\forall\in\Q$, $r\left( \prod_{p\text{prime}}|r|_p \right) =1$.
\begin{definition}
	Let $k$ be a field. A \textit{\textbf{valuation}} is a function
	\begin{align*}
		\nu: k &\longrightarrow\R_{>}  \\
	\end{align*}
	if
	\begin{enumerate}
		\item $\nu(ab)=\nu(a)+\nu(b)$
		\item $\nu(a+b)\geq \min(\nu(a),\nu(b))$
	\end{enumerate}
\end{definition}
$\nu(a+b)=\min$ iff what?


Novikov field, Stanley-Reiser varieties. What ideals correspond to these varieties? Monomial ideals This program will give me monomial ideals.

Question. Given a simplicial complex $T$ when have thus produced a monomial ideal $I_T$. So this \texttt{gfan} will construct a list of monomial ideals. What ideal should you give the program to get $I_T$? Anwser: $I_T$.

Now choose a better variety, an irreducible $X\subset \mathbb{P}^n$. It has a prime ideal $J$.

Explicit problem: among all triangulations of 3-sphere, there is a very strange one: Br\"uckner-Gr\"unbaum sphere.

Somehow this leads to a Calabi-Yau threefold of degree 20 $X^{(3)}\hookrightarrow\mathbb{P}^7$, that is new. Now you have to smooth it.

Why is this important? We actually want some $X^{(4)}\hookrightarrow\mathbb{P}^8$ Hyperk\"ahler variety of type Kun 8. We're studying so-called $\C\mathbb{P}^2_9=$ B. K\"uhne. Claim. There exists a triangulation of $\C\mathbb{P}^2$ with only 9 vertices.

\begin{thm}[who?]
	If $M^(d)$ is a manifold of dimension $d$ with $N\geq d+2$ vertices and $N\leq d+4$ then $M$ is a sphere.
\end{thm}
For example, if $M^{(4)}$ with $\leq 8$ vertices you get a sphere or $\C\mathbb{P}^2_9$. See article On K\"uhnel 9-vertex complex projective plane.

{\color{cyan}Link with any vertices is a triangulation of sphere. $\operatorname{Link}(V)=S_p^3$.}

\end{document}
